\documentclass[11pt]{article}
\usepackage[margin=1in]{geometry}
\usepackage{booktabs}
\usepackage{hyperref}
\usepackage{parskip}
\usepackage{enumitem}
\usepackage{float}

\title{\textbf{Final Project Proposal:\\
Urban Flood Detection via Cross-Modal Weak Supervision}}
\author{Romit Basak}
\date{CS5330 Pattern Recognition and Computer Vision --- Spring 2026}

\begin{document}
\maketitle

% ============================================================================
\section*{Problem Statement}
% ============================================================================

Urban flood detection from satellite imagery is an open problem. Optical satellites
produce high-quality water maps but fail during cloud cover --- precisely when floods
are most likely to occur. Synthetic Aperture Radar (SAR) penetrates clouds but produces
complex backscatter signals that are difficult to interpret in urban areas, where radar
pulses bounce off building walls rather than reflecting directly from water surfaces
(the double-bounce effect).

This project builds a \textbf{generalizable} deep learning pipeline for urban flood
detection that can be deployed in any city worldwide. The core idea: use
\textbf{clear-sky optical imagery as pseudo-ground-truth} to train a SAR-based flood
detector. When a flood event coincides with clear skies, optical water indices (NDWI)
can generate pixel-precise flood labels automatically. These labels train a U-Net
segmentation model on co-registered SAR imagery. The trained SAR model can subsequently
detect floods under cloudy conditions where optical sensors fail entirely --- making it
useful for any urban region regardless of local cloud climatology.

% ============================================================================
\section*{Data}
% ============================================================================

\textbf{Primary satellite data} will be accessed via Google Earth Engine (GEE), which
provides free, analysis-ready access to the full Sentinel-1 SAR and Sentinel-2 optical
archives:

\begin{itemize}[noitemsep]
    \item \textbf{Sentinel-1 SAR} (\texttt{COPERNICUS/S1\_GRD}): 10m resolution,
    C-band, VV and VH polarizations. Cloud-penetrating. Used as model input.
    \item \textbf{Sentinel-2 optical} (\texttt{COPERNICUS/S2\_SR\_HARMONIZED}):
    10m resolution, multi-spectral. Used only during clear-sky periods to generate
    NDWI-based pseudo-labels. Not available at inference time.
    \item \textbf{Landsat 8/9} (\texttt{LANDSAT/LC08/C02/T1\_L2}): 30m fallback
    when Sentinel-2 coverage is insufficient. Downsampled to match SAR resolution.
\end{itemize}

\textbf{Public benchmark datasets} will be used for pre-training and evaluation:

\begin{itemize}[noitemsep]
    \item \textbf{UrbanSARFloods} (Zhao et al., CVPR Workshop 2024): 8,879 chips
    across 18 flood events, Sentinel-1 intensity and InSAR coherence, urban focus.
    Primary benchmark and source of pre-labeled training chips.
    \item \textbf{Sen1Floods11} (Bonafilia et al., CVPR Workshop 2020): 4,831
    hand-labeled chips across 11 events, used for initial pretraining.
\end{itemize}

\textbf{Pseudo-label generation pipeline}: For $\sim$27 SAR--optical image pairs spanning 3 major flood events across 8 city regions (Houston Harvey 2017, Valencia Spain 2024, Brisbane 2022), the pipeline
finds coincident SAR and clear optical acquisitions (within 48 hours), computes NDWI
from the optical image, subtracts a per-city dry-season median baseline to isolate
anomalous water, and generates confidence-weighted binary flood masks. This yields
additional training pairs beyond what the public benchmarks provide, with controllable
geographic diversity.

% ============================================================================
\section*{Model}
% ============================================================================

The core model is a \textbf{U-Net} segmentation architecture, chosen for its strong
track record on satellite imagery and its encoder-decoder structure that preserves
spatial resolution --- essential for 10m-pixel flood boundary delineation.

\textbf{Core input channels} (per image chip, $256 \times 256$ pixels):
\begin{itemize}[noitemsep]
    \item VV backscatter (Sentinel-1)
    \item VH backscatter (Sentinel-1)
    \item Local variance (per-channel $5\times5$ filter, discriminates radar shadow from open water via speckle texture)
    \item Incidence angle (per-pixel, from SAR metadata)
    \item ERA5 precipitation accumulation (72-hour pre-event, optional conditioning channel)
\end{itemize}

\textbf{Training strategy}:
\begin{enumerate}[noitemsep]
    \item Pre-train on Sen1Floods11 (large, diverse, hand-labeled).
    \item Fine-tune on UrbanSARFloods + pseudo-labels from 2--3 collected flood events.
    \item Evaluate on held-out cities from UrbanSARFloods to test generalizability.
\end{enumerate}

\textbf{Loss function}: Weighted combination of binary cross-entropy and IoU loss,
with confidence weighting from the pseudo-label pipeline to down-weight uncertain
pixels near flood boundaries.

\textbf{Baselines}: Otsu thresholding and Random Forest on raw backscatter, replicating
the Al Mehedi et al.\ (2025) experimental setup for direct comparison. Primary
evaluation metric: F1-score on the flood class (accuracy is misleading due to severe
class imbalance).

% ============================================================================
\section*{Planned Extensions (Post-Submission)}
% ============================================================================

The following components are scoped out of the April 22 submission due to time
constraints, but are planned extensions for continued thesis work:

\begin{itemize}[noitemsep]
    \item \textbf{Radar shadow channel}: Building footprints (Google Open Buildings)
    combined with SAR geometry to compute per-pixel radar shadow probability, helping
    the model distinguish shadow from water. Requires non-trivial geometric computation.
    \item \textbf{L-band generalization}: Extending the pipeline to ALOS-2 L-band SAR
    to test whether the model transfers across radar wavelengths.
    \item \textbf{Seasonal baseline expansion}: Scaling pseudo-label generation to
    full wet seasons (hundreds of training pairs per city) rather than 2--3 events.
    \item \textbf{Kolkata deployment}: Applying the trained generalizable model to
    Kolkata monsoon flood detection, connecting this work to ongoing thesis research
    with Prof.\ Auroop Ganguly.
\end{itemize}

% ============================================================================
\section*{Compute}
% ============================================================================

\textbf{Data preprocessing} will run entirely on \textbf{Google Earth Engine} (free,
cloud-based), handling the petabyte-scale satellite archive without local compute.
GEE exports processed image chips directly to Google Drive.

\textbf{Model training} will use \textbf{Google Colab Pro} (available via university
access), which provides A100 GPU instances. U-Net at $256 \times 256$ input with
$\sim$30M parameters trains in approximately 2--4 hours per run. Total estimated GPU
time across all experiments: 20--30 hours.

The Khoury College GPU cluster (RTX 3090, SLURM) is available as a fallback.
No paid cloud services are required beyond existing university resources.

% ============================================================================
\section*{Connection to CS5330}
% ============================================================================

This project applies core course concepts directly to a real remote sensing problem:
semantic segmentation (U-Net), transfer learning (ImageNet $\to$ SAR domain),
cross-modal learning (optical pseudo-labels for SAR training), and physics-based
feature engineering. The generalizable flood detection system is a direct extension
of the deep network recognition and transfer learning themes developed in Projects 4--5.

\end{document}
